% AUTOGENERATED from hw1_math_solutions.tex by build.py -- do not edit; edit the _solutions file.
\documentclass{cs1840}

\usepackage{parskip}

% Q1 action labels
\newcommand{\noheal}{\texttt{Don't Heal}}
\newcommand{\heal}{\texttt{Heal}}

\hwnum{1}
\hwtitle{Markov Decision Processes}
\hwsubtitle{Math}
\hwdue{September 28}

\begin{document}
\makehwtitle

\hwinstructions

Throughout this assignment, assume that:
state and action spaces are finite and nonempty;
a value at state $s$ means the expected return when initialized at $s$, even if $s$ would not be reached from another starting state;
in finite-horizon problems, a policy may depend on the time step, no actions are taken at the final time $H$, and no terminal reward is given.

\begin{problem}{Pok\'emon as an MDP}

You're up against a never-ending Elite Four, and the only Pok\'emon you have is your Charizard.
Even worse, you didn't bring any healing items.
Your Charizard's status is tracked by its damage level in $\{0, 1, 2\}$, where $0$ is full health and $2$ is one hit from fainting.
Writing $F$ for the fainted state, we have $\Scal = \{0, 1, 2, F\}$.

Since you only spam Flare Blitz, your Charizard wins every battle but gets injured in the process.
More precisely, if you battle at damage level $s$, your Charizard will faint after winning, with probability $s/2$.
Otherwise, your Charizard wins, taking only minor recoil damage, and has its damage level increase to $s+1$, with probability $1-s/2$.
(There is one exception to this below, when buying a Potion.)
At damage level $2$, fainting is therefore certain, and when your Charizard is fainted, trying to battle will keep it fainted and do nothing instead.

Fortunately, you have your credit card with you!
At each time step, you may choose either to battle without an item ($a=\noheal$) or purchase and immediately use an item ($a=\heal$).
Thus, $\Acal=\{\noheal, \heal\}$.
When your Charizard is not fainted, you can spend $p>0$ Pok\'eDollars to buy a Potion, which will bring your Charizard to full health.
Moreover, the Potion also allows you to win a battle on the \emph{same} time step without taking damage.
When your Charizard is fainted, you can spend $2p$ Pok\'eDollars to buy a Revive, which will simply bring your Charizard to full health.

Lastly, each battle you win also grants $1$ Pok\'eDollar in prize money.
Your reward at each time step is your change in Pok\'eDollars, and rewards are earned on the current time step, before moving to the next step.
Note that you still receive the prize even if Charizard faints after winning.
You may also assume that you can always buy an item at any step regardless of past earnings.

We will now model this problem as an infinite-horizon MDP $\mathcal{M} \coloneqq (\Scal, \Acal, P, r, \gamma)$, with discount factor $\gamma\in(0,1)$.

\begin{enumerate}
    \item \points{8} Write down the transition model $P(s'|s, a)$ and reward function $r(s, a)$.

\begin{s}
    The transition model $P(s'|s, a)$ and reward function $r(s, a)$ are

\begin{center}
\begin{tabular}{|c|c|c|c|c|c|c|}
    \hline
    $s$ & $a$ & $P(0 \mid s,a)$ & $P(1 \mid s,a)$ & $P(2 \mid s,a)$ & $P(F \mid s,a)$ & $r(s,a)$ \\
    \hline
    $0$ & \noheal & $0$ & $1$ & $0$ & $0$ & $1$ \\
    $0$ & \heal   & $1$ & $0$ & $0$ & $0$ & $1-p$ \\
    \hline
    $1$ & \noheal & $0$ & $0$ & $1/2$ & $1/2$ & $1$ \\
    $1$ & \heal   & $1$ & $0$ & $0$ & $0$ & $1-p$ \\
    \hline
    $2$ & \noheal & $0$ & $0$ & $0$ & $1$ & $1$ \\
    $2$ & \heal   & $1$ & $0$ & $0$ & $0$ & $1-p$ \\
    \hline
    $F$ & \noheal & $0$ & $0$ & $0$ & $1$ & $0$ \\
    $F$ & \heal   & $1$ & $0$ & $0$ & $0$ & $-2p$ \\
    \hline
\end{tabular}
\end{center}

\end{s}

    \item \points{10} Consider the threshold policy $\pi_C$ for $C\in\{0,1,2\}$, given by
    \[
        \pi_C(s) =
        \begin{cases}
            \heal,   & s\geq C ~\text{or}~s=F \\
            \noheal, & \text{otherwise}
        \end{cases}.
    \]
Derive an equation for $V^{\pi_C}(F)$ and $V^{\pi_C}(s)$ for numeric $s\geq C$ in terms of $V^{\pi_C}(0)$.
Then, derive a recursive equation for $V^{\pi_C}(s)$ for numeric $s<C$ in terms of $V^{\pi_C}(0)$ and $V^{\pi_C}(s+1)$.

    \textit{Hint:} Use the Bellman consistency equations.

\begin{s}
    We have three cases to consider, each one applying the Bellman consistency equation.

\textbf{Case 1 ($s = F$):} The policy is \heal, the next state is 0 with probability 1, and the reward is $-2p$:

\[
V^{\pi_C}(F) = -2p + \gamma V^{\pi_C}(0)
\]

\textbf{Case 2 ($s \geq C$):} The policy is \heal, the next state is 0 with probability 1, and the reward is $1 - p$:

\[
V^{\pi_C}(s) = 1 - p + \gamma V^{\pi_C}(0)
\]

\textbf{Case 3 ($s < C$):} The policy is \noheal, the next state is $s + 1$ with probability $1 - s/2$ and $F$ with probability $s / 2$, and the reward is $1$:

\[
V^{\pi_C}(s) = 1 + \gamma \left( \left(1 - \frac{s}{2} \right) V^{\pi_C}(s + 1) + \frac{s}{2} V^{\pi_C}(F) \right)
\]

Substitute in $V^{\pi_C}(F)$ from Case 1:

\[
V^{\pi_C}(s) = 1 + \gamma \left[ \left(1 - \frac{s}{2} \right) V^{\pi_C}(s + 1) + \frac{s}{2} (-2p+ \gamma V^{\pi_C}(0)) \right]
\]
\end{s}

    \item \points{16} For $\gamma=1/2$ and $p=0.25$, find the value of $C^\star$ such that $V^{\pi_{C^\star}}(s) \geq V^{\pi_C}(s)$ for all $s\in\Scal$ and $C\in\{0,1,2\}$.


\begin{s}
    We can test the threshold policy for each value of $C$ with $\gamma = 1/2$ and $p = 0.25$:

\textbf{C = 0:} The policy is always \heal: 

\[
V^{\pi_C}(0) = \frac{3}{4} + \frac{1}{2} V^{\pi_C}(0) \Rightarrow V^{\pi_C}(0) = \frac{3}{2}
\]

Similarly, $V^{\pi_C}(1) = V^{\pi_C}(2) = 3/2$. For $s = F$,

\[
V^{\pi_C}(F) = -\frac{1}{2} + \frac{1}{2}V^{\pi_C}(0) = \frac{1}{4}
\]

\textbf{C = 1:} The policy is \noheal on 0 and \heal otherwise:

\begin{align*}
V^{\pi_C}(0) &= 1 + \frac{1}{2}V^{\pi_C}(1) \\
V^{\pi_C}(1) &= V^{\pi_C}(2) = \frac{3}{4} + \frac{1}{2}V^{\pi_C}(0)
\end{align*}

so $V^{\pi_C}(0) = 11/6$ and $V^{\pi_C}(1) = 5/3$. For $s = F$,

\[
V^{\pi_C}(F) = -\frac{1}{2} + \frac{1}{2}V^{\pi_C}(0) = \frac{5}{12}
\]

\textbf{C = 2:} The policy is \noheal on 0, 1 and \heal on 2, F:

\begin{align*}
V^{\pi_C}(0) &= 1 + \frac{1}{2}  V^{\pi_C}(1) \\
V^{\pi_C}(1) &= 1 + \frac{1}{2} \left[ \frac{1}{2}V^{\pi_C}(2) + \frac{1}{2} V^{\pi_C}(F) \right] \\
V^{\pi_C}(2) &= \frac{3}{4} + \frac{1}{2} V^{\pi_C}(0) \\
V^{\pi_C}(F) &= -\frac{1}{2}+ \frac{1}{2} V^{\pi_C}(0)
\end{align*}

Solving gives $V^{\pi_C}(0) = 7/4$, $V^{\pi_C}(1) = 3/2$, $V^{\pi_C}(2) = 13/8$, and $V^{\pi_C}(F) = 3/8$.

The results from each of the cases can be summarized in the following table. The highest scores in each column come from the row $C = 1$, so $C^* = 1$:

\[
\begin{array}{c|c|c|c|c|}
C & V^{\pi_C}(0) & V^{\pi_C}(1) & V^{\pi_C}(2) & V^{\pi_C}(F) \\ \hline
0 & \frac{3}{2} & \frac{3}{2} & \frac{3}{2} & \frac{1}{4} \\
1 & \frac{11}{6} & \frac{5}{3} & \frac{5}{3} & \frac{5}{12} \\
2 & \frac{7}{4} & \frac{3}{2} & \frac{13}{8} & \frac{3}{8}
\end{array}
\]


\end{s}

    \item \points{16} Keeping $\gamma=1/2$ and $p=0.25$, show that $\pi_{C^\star}$ is optimal among all policies, not just the three threshold policies.

    \textit{Hint:} Use the Bellman optimality equations.

\begin{s}
    The Bellman optimality equation states

    \[
    V^*(s) = \max_a [r(s, a) + \gamma \mathbb{E}_{s' \sim P(s, a)} V^*(s')]
    \]

    for a unique value function. We can check each state and action pair to see if $V^{\pi_1} = (11/6, 5/3, 5/3, 5/12)$ satisfies the equation.

    \textbf{s = 0:}

\[
\max_a [r(0, a) + \gamma \mathbb{E}_{s' \sim P(0, a)} V^{\pi_1}(s')] = \max \left[ 1 - \frac{1}{4} + \frac{1}{2}V^{\pi_1}(0), 1 + \frac{1}{2}V^{\pi_1}(1)\right]
\]
\[
= \max \left[ \frac{5}{3}, \frac{11}{6}\right] =  V^{\pi_1}(0)
\]

    \textbf{s = 1:}

\[
\max_a [r(1, a) + \gamma \mathbb{E}_{s' \sim P(1, a)} V^{\pi_1}(s')]
= \max \left[
1 - \frac{1}{4} + \frac{1}{2}V^{\pi_1}(0),
1 + \frac{1}{2}\left(\frac{1}{2}V^{\pi_1}(2) + \frac{1}{2}V^{\pi_1}(F)\right)
\right]
\]
\[
= \max \left[
\frac{5}{3}, \frac{73}{48}
\right]
= V^{\pi_1}(1)
\]

    \textbf{s = 2:}

\[
\max_a [r(2, a) + \gamma \mathbb{E}_{s' \sim P(2, a)} V^{\pi_1}(s')]
= \max \left[
1 - \frac{1}{4} + \frac{1}{2}V^{\pi_1}(0),
1 + \frac{1}{2}V^{\pi_1}(F)
\right]
\]

\[
= \max \left[
\frac{5}{3}, \frac{29}{24}
\right]
= V^{\pi_1}(2)
\]

    \textbf{s = F:}

\[
\max_a [r(F, a) + \gamma \mathbb{E}_{s' \sim P(F, a)} V^{\pi_1}(s')]
= \max \left[
-2\left(\frac{1}{4}\right) + \frac{1}{2}V^{\pi_1}(0),
\frac{1}{2}V^{\pi_1}(F)
\right]
\]

\[
= \max \left[
\frac{5}{12}, \frac{5}{24}
\right]
= V^{\pi_1}(F)
\]

Thus, $V^{\pi_1}(s)$ satisfies the Bellman optimality equations, demonstrating that $\pi_1$ is the optimal policy.
 
\end{s}
\end{enumerate}

\end{problem}

\newpage


\begin{problem}{Bellman optimality for \texorpdfstring{$Q^\star_h$}{Q\*\_h}}

Consider an (undiscounted) finite-horizon MDP, with horizon $H$.
Policies are (possibly) time-dependent, given by tuples $\pi=(\pi_0,\ldots,\pi_{H-1})$.
In lecture, we argued (without proof) that there exists a deterministic policy $\pi^\star$ such that
\[
    V^{\pi^\star}_h(s) = \max_{\pi} V^\pi_h(s)
\]
for all $s\in\Scal$ and $h\in\{0,\dots,H-1\}$.
In other words, there exists a single policy that is simultaneously optimal at all states and time steps.
(Note that we use the terminal conventions $V^\pi_H(s)=0$ and $Q^\pi_H(s,a)=0$ for \emph{all} policies $\pi$.)
From the Bellman consistency, this also implies
\[
    Q^{\pi^\star}_h(s,a) = \max_\pi Q^\pi_h(s,a)
\]
for all $s\in\Scal$, $a\in\Acal$, and $h\in\{0,\dots,H-1\}$.
We can thus write $V^\star_h \coloneqq V^{\pi^\star}_h$ and $Q^\star_h \coloneqq Q^{\pi^\star}_h$ to refer to the value functions of the optimal policy.
Now, recall from lecture that
\[
    V^\star_h(s) = \max_a \left\{ r(s, a) + \mathbb{E}_{s' \sim P(\cdot | s, a)}\left[V^\star_{h+1}(s')\right] \right\}
\]
for all $s\in\Scal$ and $h\in\{0,\dots,H-1\}$.
These are the Bellman optimality equations for $V^\star_h$.
In this problem, you will prove that
\[
    Q^\star_h(s,a) = r(s,a) + \mathbb{E}_{s'\sim P(\cdot|s,a)} \left[\max_{a'} Q^\star_{h+1}(s',a') \right]
\]
for all $s\in\Scal$, $a\in\Acal$, and $h\in\{0,\dots,H-1\}$.

\begin{enumerate}
    \item \points{5} Prove that $Q^\star_h (s, \pi^\star_h(s)) = \max_a Q^\star_h (s,a)$.

    \textit{Hint:} For a deterministic policy $\pi$, recall that $Q^\pi_h(s,\pi_h(s))=V^\pi_h(s)$.
    You might also find the Bellman optimality equations and the Bellman consistency equations useful.

\begin{s}
Since $\pi^*$ is deterministic,
\[
    Q^*_h\big(s,\pi^*_h(s)\big) = Q^{\pi^*}_h\big(s,\pi^*_h(s)\big) = V^{\pi^*}_h(s) = V^*_h(s)
\]
By the relationship between $V^\pi$ and $Q^\pi$,
\[
    Q^*_h(s,a) = r(s,a) + \mathbb{E}_{s' \sim P(\cdot \mid s,a)}\left[V^*_{h+1}(s')\right]
\]
Substituting that into the Bellman optimality equation for $V^*_h$,
\[
    V^*_h(s) = \max_a \left[ r(s,a) + \mathbb{E}_{s' \sim P(\cdot \mid s,a)} V^*_{h+1}(s') \right] = \max_a Q^*_h(s,a)
\]
Combine the two results:
\[
    Q^*_h\big(s,\pi^*_h(s)\big) = V^*_h(s) = \max_a Q^*_h(s,a)
\]

\end{s}

    \item \points{5} Prove the Bellman optimality for $Q^\star_h$.

    \textit{Hint:} For time steps $h=0,\dots,H-2$, start from the Bellman consistency equations, and use your result from (a).
    You can handle $h=H-1$ directly using the terminal convention.

\begin{s}
The relationship between $Q$ and $V$ is defined as $Q^\pi_h(s, a) = r(s, a) + \mathbb{E}_{s' \sim P(\cdot \mid s, a)} \left[ V^\pi_{h + 1}(s')\right]$. Let $\pi^*$ be the optimal policy. For $h \in \{0, \dots, H-2\}$,

\begin{align*}
    Q^*_h(s, a) &= Q^{\pi^*}_h(s, a) \\
    &= r(s, a) + \mathbb{E}_{s' \sim P(\cdot \mid s, a)} \left[ V^{\pi^*}_{h + 1}(s') \right] \\
    &= r(s, a) + \mathbb{E}_{s' \sim P(\cdot \mid s, a)} \left[ V^{*}_{h + 1}(s') \right]
\end{align*}

Since $h + 1 \leq H - 1$, part (a) applies at step $h + 1$, so $V^*_{h + 1}(s') = \max_{a'}Q^*_{h+1}(s', a')$. Thus,

\[
Q^*_h(s, a) = r(s, a) + \mathbb{E}_{s' \sim P(\cdot \mid s, a)} \left[ \max_{a'} Q^*_{h + 1}(s', a') \right]
\]

For $h = H - 1$, $V^*_H(s') = 0$ and $Q^*_H(s', a') = 0$, so both sides of the equation equal $r(s, a)$:

\[
Q^*_{H - 1}(s, a) = r(s, a) = r(s, a) + \mathbb{E}_{s' \sim P(\cdot \mid s, a)} \left[ \max_{a'} Q^*_{H}(s', a') \right]
\]
\end{s}


\end{enumerate}

\end{problem}

\newpage

\begin{problem}{Finitely approximating an infinite horizon MDP}

Recall that a finite-horizon MDP is given by a tuple $\mathcal{M}_H \coloneqq (\Scal, \Acal, P, r, H)$, where returns are (undiscounted) sums of the specified rewards $r$.
(We'll leave out the initial distribution in this problem.)
The reward function is $r: \Scal\times \Acal \rightarrow [0,1]$, where $r(s,a)$ is the instantaneous reward at state-action pair $(s,a)$, independent of the time step.
Now, consider modifying the MDP to obtain $\overline{\mathcal{M}}_H \coloneqq (\Scal, \Acal, P, \overline{r}, H)$, where the reward function now depends on the time step.
More precisely, the reward function is $\overline{r}:\Scal\times\Acal \times\{0, \ldots, H-1\} \rightarrow [0,1]$, where $\overline{r}(s,a,h)$ is the instantaneous reward at time step $h$.
We now define
\[
    \overline{V}^{\pi}_h(s) \coloneqq \mathbb{E}^{\pi}\left[ \sum_{t = h}^{H-1} \overline{r}(s_t, a_t,t) \Big| s_h = s \right].
\]

\begin{enumerate}
    \item \points{6} Describe a dynamic programming algorithm to compute
    both $\overline{V}^{\star}_h(s)$ and $\overline{\pi}^{\star}_h(s)$, the optimal value function and policy in $\overline{\mathcal{M}}_H$.
    (You do not need to prove your program is correct, though it should be clear how it generalizes the original algorithm.)
\end{enumerate}

\begin{s} We can use backward dynamic programming to solve for $\overline{V}^*_h(s)$ and $\overline{\pi}^*_h(s)$. At the end of the finite horizon,

\[
V_H^*(s)=0
\]

For $h=H-1,\ldots,0$,

\[
Q_h^*(s,a)
=
\overline r(s,a,h)
+
\mathbb{E}_{s' \sim P(s,a)}
\left[V_{h+1}^*(s')\right]
\]

and

\[
V_h^*(s)
=
\max_a Q_h^*(s,a)
\]

Thus, the optimal policy is

\[
\pi_h^*(s)
=
\arg\max_a Q_h^*(s,a)
\]

\end{s}

Now, consider an infinite-horizon discounted MDP, $\mathcal{M}_\gamma \coloneqq (\Scal, \Acal, P, r, \gamma)$, where $r: \Scal\times\Acal \rightarrow [0,1]$ is the usual reward function.
You will define a finite-horizon MDP $\overline{\mathcal{M}}_{\overline H} \coloneqq (\Scal, \Acal, \overline P, \overline r, \overline H)$ which, in a precise sense, approximates $\mathcal{M}_\gamma = (\Scal, \Acal, P, r, \gamma)$.
To start, we have specified that both MDPs share the same state and action spaces.
It remains for you to specify $\overline P$, $\overline r$ and $\overline H$.

Let $\overline{V}^{\pi}_0(s)$ refer to the value function for the finite-horizon MDP $\overline{\mathcal{M}}_{\overline H}$, and let $V^{\pi}_{\gamma}(s)$ be the value function for the infinite-horizon discounted MDP $\mathcal{M}_\gamma$.
Fix $\eps\in(0,1)$.
Our goal is the following: for all $s\in\Scal$, and for all stationary policies $\pi$, we want
\[
    \left| V^{\pi}_{\gamma }(s) - \overline{V}^{\pi}_0(s) \right| \leq \eps.
\]
(Recall that a stationary policy $\pi$ does not depend on the time step $h$.
We restrict to these so that the same $\pi$ can be used for both the finite and infinite horizons.)

\begin{enumerate}
    \setcounter{enumi}{1}
    \item \points{12}
    We will first establish an analytic fact.
    We want to find an integer $H_\eps$ that has the following property.
    For any sequence $x_0, x_1, \dots$ of numbers such that $x_t\in[0,1]$,
    \[
        \sum_{t=0}^\infty \gamma^t x_t - \sum_{t=0}^{H_\eps-1} \gamma^t x_t \leq \eps.
    \]
    (Recall that $\gamma\in(0,1)$.)
    Show that it suffices to take
    \[
        H_\eps \coloneqq \left\lceil \frac{\ln\big[1/((1-\gamma)\eps)\big]}{1-\gamma} \right\rceil.
    \]
    \textit{Hint:} A useful fact is that $-\ln(x) \geq 1-x$ for all $x>0$.
    (This follows from $1+x \leq \exp(x)$.)
\end{enumerate}

\begin{s} Since $x_t \in [0,1]$,

\[
\sum_{t=0}^{\infty} \gamma^t x_t
-
\sum_{t=0}^{H_\epsilon-1} \gamma^t x_t
=
\sum_{t=H_\epsilon}^{\infty} \gamma^t x_t
\leq
\sum_{t=H_\epsilon}^{\infty} \gamma^t
=
\frac{\gamma^{H_\epsilon}}{1-\gamma}
\]

Using the inequality from the hint, $-\ln \gamma \ge 1-\gamma$, so $\ln \gamma \le -(1-\gamma)$. Therefore,

\[
\gamma^{H_\epsilon}
=
e^{H_\epsilon \ln \gamma}
\leq
e^{-H_\epsilon(1-\gamma)}
\]

By the definition of $H_\epsilon$,

\[
H_\epsilon
\geq
\frac{\ln\big[1/((1-\gamma)\epsilon)\big]}{1-\gamma}
\]

so

\[
\gamma^{H_\epsilon}
\leq
e^{-\ln \left(\frac{1}{(1-\gamma)\epsilon}\right)}
=
(1-\gamma)\epsilon.
\]

Therefore, we have

\[
\frac{\gamma^{H_\epsilon}}{1-\gamma}
\leq
\epsilon
\]

so

\[
\sum_{t=H_\epsilon}^{\infty} \gamma^t x_t
\leq \epsilon
\]
    
\end{s}

\begin{enumerate}
    \setcounter{enumi}{2}
    \item \points{6}
    Specify $\overline P$, $\overline r$ and $\overline H$.
    These should be written only in terms of $\eps$ and any other relevant quantities in the tuple $\mathcal{M}_\gamma$.
    Use the horizon bound from part (b).

    \textit{Hint:} $\overline r$ is allowed to be time-dependent.

\begin{s}
    We can choose

    \[
    \overline{P}(s'|s, a) = P(s'|s, a) \qquad \overline{r}(s, a, h) = \gamma^hr(s, a) \qquad \overline{H} = H_\epsilon
    \]

    to match the horizon bound from part (b).
\end{s}

    \item \points{6} Briefly argue why your choices imply our desired condition.
    In particular, explain what motivated your choice of $\overline P$.

\begin{s}
    Since $\overline{P} = P$, the trajectory distributions under the same policy $\pi$ are identical in both MDPs. The two value functions are

\[
\overline{V}_0^\pi(s)
=
\mathbb{E}^\pi
\left[
\sum_{t=0}^{H_\epsilon-1}
\gamma^t r(s_t,a_t)
\,\middle|\, s_0=s
\right]
\]

\[
V_\gamma^\pi(s)
=
\mathbb{E}^\pi
\left[
\sum_{t=0}^{\infty}
\gamma^t r(s_t,a_t)
\,\middle|\, s_0=s
\right]
\]

so


\[
\begin{aligned}
V_\gamma^\pi(s) - \overline{V}_0^\pi(s)
&=
\mathbb{E}^\pi
\left[
\sum_{t=H_\epsilon}^{\infty}
\gamma^t r(s_t,a_t)
\,\middle|\, s_0=s
\right] \\
&\leq
\mathbb{E}^\pi
\left[
\sum_{t=H_\epsilon}^{\infty}
\gamma^t
\,\middle|\, s_0=s
\right] \\
&=
\sum_{t=H_\epsilon}^{\infty}\gamma^t \\
&\leq \epsilon
\end{aligned}
\]

where the last inequality follows from part (b) and the fact that $r(s_t, a_t) \in [0, 1]$. Note that $V_\gamma^\pi(s) - \overline{V}_0^\pi(s)$ is also nonnegative, since $r \geq 0$. Therefore, for every state $s$ and stationary policy $\pi$,

\[
\left|V_\gamma^\pi(s) - \overline{V}_0^\pi(s)\right|
\leq \epsilon
\]

\end{s}
\end{enumerate}

Because of the above, we naturally refer to $\overline H$ as the $\eps$-\emph{effective horizon}.
The interpretation is that we can approximate any infinite-horizon discounted MDP by modifying a finite-horizon MDP with sufficiently long horizon.
This turns out to be quite useful; by running your program in (a) on the MDP $\overline{\mathcal{M}}_{\overline H}$ you constructed in (c), we get an $\eps$-optimal value function in the \emph{infinite-horizon} MDP!
This approach leads to a nearly identical sequence of value estimates to value iteration, differing only by a scaling factor at each iteration.
(There is no graded question here, but we encourage you to think about why this is true!)

\end{problem}

\end{document}
